
The details of the event generation for the $m_W$ measurement differ between measurements due to 
the assumed $W$-boson width ${\mathrm \Gamma}_W$ and to a restriction on the generated $\ell\nu$ invariant 
mass range in the CDF sample.  These lead to $\delta m_W$ corrections on the direct fits to these 
samples.  

The assumed ${\mathrm \Gamma}_W$ values used by the experiments are as follows: 2089.5 MeV for the CDF 
measurement; 2099 and 2100 MeV for the D0 measurements with 1.0~fb$^{-1}$ and 4.3~fb$^{-1}$, 
respectively; 2094 MeV for the ATLAS measurement; and 2085 for the LHCb measurement.  Using the  
SM prediction of ${\mathrm \Gamma}_W = 2089.5 \pm 0.6$~MeV leads to $\delta m_W^{\mathrm \Gamma}=0.0$, 1.4,
1.5, 0.7, and $-0.7$~MeV for the CDF, D0 1.0 fb$^{-1}$, D0 4.3 fb$^{-1}$, ATLAS, and LHCb measurements,
respectively.

The {\textsc{ResBos-C}} events used in the fit to the CDF data includes an $\ell\nu$ invariant mass 
requirement of $m_{\ell\nu}<150$~GeV.  Differences of up to 10\% are observed between {\textsc{Resbos-C}} 
and {\textsc{ResBos2}} for an invariant mass below 70~GeV, though these have a negligible effect on the 
measurement.  Using the {\textsc{ResBos2}} invariant mass distribution without any requirements leads 
to $\delta m_W^{\scriptsize\textrm{gen}} = -1.6, -3.4,$ and $-3.2$~MeV for the \mT, \pTl, and \met~distributions,
respectively, for the CDF fit results.  The measured CDF $m_W$ accounts for these effects as part of the update
of the PDF set to NNPDF3.1.  Smaller differences are observed between {\textsc{Resbos-CP}} and {\textsc{ResBos2}}, 
and there is no significant $\delta m_W^{\scriptsize\textrm{gen}}$ from the invariant mass modelling for the D0
measurement.
